\documentclass[10pt,a4paper]{article}
\usepackage[utf8]{inputenc}
\usepackage[T1]{fontenc}
\usepackage[french]{babel}
\usepackage{lmodern}
\usepackage[a4paper,margin=1.8cm,headheight=15pt]{geometry}
\usepackage{xcolor}
\usepackage{hyperref}
\usepackage{fancyhdr}
\usepackage{enumitem}
\definecolor{forest}{HTML}{2B5142}
\definecolor{muted}{HTML}{637167}
\hypersetup{colorlinks=true,linkcolor=forest,urlcolor=forest,pdftitle={TaskFlow - dossiers et fichiers}}
\pagestyle{fancy}\fancyhf{}
\fancyhead[L]{\small\textcolor{forest}{TASKFLOW / GUIDE DU PROJET}}
\fancyhead[R]{\small Architecture et données}
\fancyfoot[L]{\small Inventaire local du dépôt}
\fancyfoot[R]{\small\thepage}
\setlength{\parindent}{0pt}
\setlength{\parskip}{0.35em}
\setlist{nosep,leftmargin=1.4em}
\begin{document}
{\Huge\bfseries\textcolor{forest}{TaskFlow}}\\[0.25em]
{\Large Dossiers, fichiers et stockage des données}\\
\textcolor{muted}{Guide détaillé du dépôt local - inventaire du 8 octobre 2026}
\section*{Où sont les données ?}
TaskFlow conserve ses données métier dans MongoDB. En développement, la commande \texttt{npm run db:local} démarre MongoDB WiredTiger et écrit sur disque dans \texttt{.local/mongodb/}, à la racine du dépôt. Le chemin est construit dans \texttt{scripts/local-mongo.js}. Le dossier est caché (nom commençant par un point) et ignoré par Git dans \texttt{.gitignore}; les enregistrements ne sont donc pas commités ni envoyés sur GitHub. Les fichiers \texttt{.wt}, journaux et catalogues sont binaires : les ouvrir dans le Finder est possible, mais leur contenu ne se lit pas comme du texte et il ne faut pas les modifier manuellement.

Pour l'inspecter dans le Finder depuis la racine : \texttt{open .local/mongodb}. Pour voir les données lisiblement, utiliser l'interface/API ou les commandes d'administration locale de \texttt{docs/DONNEES.md}. L'API ne permet pas de lister tous les comptes; la commande diagnostic locale peut afficher les emails, et doit rester réservée à la base de développement. Le mot de passe est stocké sous forme de hash bcrypt. Avec Docker, la persistance est dans le volume \texttt{taskflow-data} configuré dans \texttt{compose.yaml}. Les tests utilisent des bases temporaires séparées.

\section*{Dossiers du projet}
Les dossiers générés/ignorés (dépendances, build, caches, données, résultats de test) sont décrits pour expliquer leur usage, mais leurs fichiers internes ne sont pas déroulés un par un. L'inventaire fichier par fichier ci-dessous couvre les fichiers suivis par Git, la documentation Postman locale et les PDF/journaux présents dans \texttt{output/pdf}. Le fichier local \texttt{backend/.env} est signalé sans en révéler les secrets.
\paragraph{\texttt{.github/workflows/}} Définition de la CI GitHub Actions.
\paragraph{\texttt{.local/mongodb/}} Base locale persistante, fichiers binaires WiredTiger ignorés par Git.
\paragraph{\texttt{.local/oral-preview/}} Rendus temporaires de pages pour contrôler le document oral.
\paragraph{\texttt{.local/postman-preview/}} Rendus temporaires de pages pour vérifier le guide Postman.
\paragraph{\texttt{.local/screenshots/}} Captures d'écran générées par les tests navigateur.
\paragraph{\texttt{backend/src/config/}} Configuration et connexion base de données.
\paragraph{\texttt{backend/src/controllers/}} Couche HTTP des contrôleurs.
\paragraph{\texttt{backend/src/middleware/}} Authentification et gestion des erreurs Express.
\paragraph{\texttt{backend/src/models/}} Schémas Mongoose et documents métier.
\paragraph{\texttt{backend/src/routes/}} Déclaration des routes API.
\paragraph{\texttt{backend/src/services/}} Logique métier et accès aux modèles.
\paragraph{\texttt{backend/src/utils/}} Validation et fonctions d'erreurs partagées.
\paragraph{\texttt{backend/test/}} Tests du backend.
\paragraph{\texttt{docs/oral/}} Documents de préparation de l'oral.
\paragraph{\texttt{docs/postman/}} Collection et guides des requêtes API.
\paragraph{\texttt{e2e/}} Tests fonctionnels exécutés dans un navigateur.
\paragraph{\texttt{frontend/dist/}} Build frontend généré, servi par Express en production locale.
\paragraph{\texttt{frontend/src/components/}} Composants React réutilisables.
\paragraph{\texttt{frontend/src/pages/}} Pages de l'application.
\paragraph{\texttt{output/pdf/}} PDF livrables et journaux de compilation.
\paragraph{\texttt{scripts/}} Outils de configuration et de vérification locale.
\paragraph{\texttt{test-results/}} Résultats générés par Playwright.
\paragraph{\texttt{node\_modules/}} Dépendances téléchargées par npm ; contenu généré, non détaillé fichier par fichier.
\section*{Fichiers, un par un}
\subsection*{.github}
\paragraph{\texttt{.github/workflows/ci.yml}} Décrit les vérifications automatisées lancées par GitHub Actions.
\subsection*{Fichiers à la racine}
\paragraph{\texttt{.gitignore}} Liste les fichiers locaux ou générés à ne pas versionner, notamment .local, .env, node\_modules et dist.
\paragraph{\texttt{README.md}} Guide principal : démarrage, architecture, routes API, sécurité, bonus et soutenance.
\subsection*{backend}
\paragraph{\texttt{backend/.env.example}} Modèle des variables de configuration du serveur, sans secret réel.
\paragraph{\texttt{backend/package.json}} Dépendances, scripts, configuration Jest et métadonnées du backend.
\paragraph{\texttt{backend/src/app.js}} Construit Express, branche les routes, expose santé/OpenAPI/Swagger et sert le frontend compilé.
\paragraph{\texttt{backend/src/config/db.js}} Connexion et déconnexion Mongoose.
\paragraph{\texttt{backend/src/config/env.js}} Charge et expose PORT, MONGODB\_URI et JWT\_SECRET.
\paragraph{\texttt{backend/src/controllers/authController.js}} Transforme inscription, connexion et changement de mot de passe en réponses HTTP.
\paragraph{\texttt{backend/src/controllers/taskController.js}} Gère les réponses HTTP des opérations de liste, création, lecture, modification et suppression.
\paragraph{\texttt{backend/src/middleware/auth.js}} Vérifie le Bearer JWT et l'existence du compte avant une route privée.
\paragraph{\texttt{backend/src/middleware/errors.js}} Convertit les erreurs applicatives en réponses JSON normalisées.
\paragraph{\texttt{backend/src/models/Task.js}} Schéma MongoDB des tâches et forme publique retournée par l'API.
\paragraph{\texttt{backend/src/models/User.js}} Schéma MongoDB des utilisateurs ; le hash de mot de passe est caché des lectures ordinaires.
\paragraph{\texttt{backend/src/routes/authRoutes.js}} Déclare les routes register, login et changement de mot de passe.
\paragraph{\texttt{backend/src/routes/taskRoutes.js}} Déclare les routes CRUD de tâches et impose l'authentification.
\paragraph{\texttt{backend/src/server.js}} Valide la configuration, connecte MongoDB, démarre Express et gère l'arrêt propre.
\paragraph{\texttt{backend/src/services/authService.js}} Inscrit et authentifie les comptes, hache les mots de passe et crée les JWT.
\paragraph{\texttt{backend/src/services/taskService.js}} Valide et stocke les tâches, applique filtres/tri et isole les données par propriétaire.
\paragraph{\texttt{backend/src/utils/errors.js}} Fabrique les erreurs HTTP standard de l'API.
\paragraph{\texttt{backend/src/utils/validation.js}} Validation stricte des données d'authentification, de tâche, d'identifiant et de filtres.
\paragraph{\texttt{backend/test/app.test.js}} Tests d'intégration de l'API : validation, authentification, CRUD, isolation et filtres.
\paragraph{\texttt{backend/test/env.js}} Prépare les variables nécessaires à l'environnement de test.
\subsection*{Fichiers à la racine}
\paragraph{\texttt{compose.yaml}} Configuration Docker Compose de la base MongoDB avec volume persistant.
\subsection*{docs}
\paragraph{\texttt{docs/CONSIGNES\_PROF.md}} Comparaison des consignes pédagogiques avec les choix et fonctionnalités du projet.
\paragraph{\texttt{docs/SOUTENANCE.md}} Conseils et déroulé pour préparer la démonstration orale.
\paragraph{\texttt{docs/openapi.json}} Contrat OpenAPI consommé par Swagger et décrivant routes, paramètres et réponses.
\paragraph{\texttt{docs/oral/DISCOURS\_10\_MINUTES.md}} Texte détaillé proposé pour une présentation orale d'environ dix minutes.
\paragraph{\texttt{docs/oral/oral\_taskflow.tex}} Source LaTeX du document de préparation à l'oral.
\paragraph{\texttt{docs/postman/TaskFlow.postman\_collection.json}} Collection Postman importable avec les scénarios de requêtes et scripts de test.
\paragraph{\texttt{docs/postman/TaskFlow\_Postman\_requetes.tex}} Source LaTeX du guide complet des requêtes Postman.
\paragraph{\texttt{docs/postman/TaskFlow\_Postman\_requetes.txt}} Guide texte des requêtes API et séquence de démonstration.
\subsection*{e2e}
\paragraph{\texttt{e2e/taskflow.spec.js}} Parcours de bout en bout du site dans un navigateur avec Playwright.
\subsection*{Fichiers à la racine}
\paragraph{\texttt{eslint.config.js}} Règles de lint JavaScript et React du dépôt.
\subsection*{frontend}
\paragraph{\texttt{frontend/index.html}} Document HTML hôte chargé par Vite avant le code React.
\paragraph{\texttt{frontend/package.json}} Dépendances et scripts de développement/build du frontend.
\paragraph{\texttt{frontend/src/App.jsx}} Routage React, session utilisateur et accès aux pages.
\paragraph{\texttt{frontend/src/api.js}} Fonctions HTTP du navigateur, ajout du jeton et traitement des erreurs API.
\paragraph{\texttt{frontend/src/components/Account.jsx}} Profil du compte connecté et formulaire de changement de mot de passe.
\paragraph{\texttt{frontend/src/components/AccountMenu.jsx}} Menu utilisateur avec accès au compte et déconnexion.
\paragraph{\texttt{frontend/src/components/AuthLayout.jsx}} Mise en page partagée des pages publiques de connexion et inscription.
\paragraph{\texttt{frontend/src/components/Dashboard.jsx}} Coordonne chargement des tâches, sélection, création, modification et suppression.
\paragraph{\texttt{frontend/src/components/Navbar.jsx}} Navigation principale du site.
\paragraph{\texttt{frontend/src/components/TaskFilters.jsx}} Filtres de statut, priorité et échéance, ainsi que contrôles de tri.
\paragraph{\texttt{frontend/src/components/TaskForm.jsx}} Formulaire de création et modification d'une tâche.
\paragraph{\texttt{frontend/src/components/TaskItem.jsx}} Affiche une tâche et ses actions.
\paragraph{\texttt{frontend/src/components/TaskList.jsx}} Affiche la liste de tâches et son état vide.
\paragraph{\texttt{frontend/src/index.css}} Styles globaux et présentation responsive.
\paragraph{\texttt{frontend/src/main.jsx}} Point d'entrée React ; monte l'application dans la page HTML.
\paragraph{\texttt{frontend/src/pages/LoginPage.jsx}} Formulaire et flux de connexion.
\paragraph{\texttt{frontend/src/pages/Registerpage.jsx}} Formulaire et flux de création de compte.
\paragraph{\texttt{frontend/src/taskUtils.js}} Libellés, formatage des dates et fonctions utilitaires d'affichage des tâches.
\paragraph{\texttt{frontend/vite.config.js}} Configuration Vite, serveur frontend et proxy des appels /api vers Express.
\subsection*{output}
\paragraph{\texttt{output/pdf/TaskFlow\_Postman\_requetes.log}} Journal de compilation LaTeX du guide Postman.
\paragraph{\texttt{output/pdf/TaskFlow\_Postman\_requetes.pdf}} PDF compilé du guide des requêtes Postman.
\paragraph{\texttt{output/pdf/oral\_taskflow.log}} Journal de compilation LaTeX du PDF oral.
\paragraph{\texttt{output/pdf/oral\_taskflow.pdf}} PDF compilé du document de préparation à la soutenance.
\subsection*{Fichiers à la racine}
\paragraph{\texttt{package-lock.json}} Verrouille les versions exactes des dépendances npm à installer.
\paragraph{\texttt{package.json}} Scripts racine, workspaces npm et outils partagés du projet.
\paragraph{\texttt{playwright.config.js}} Configuration des navigateurs, du serveur et des suites Playwright.
\subsection*{scripts}
\paragraph{\texttt{scripts/check-persistence.js}} Vérifie qu'une tâche survit à l'arrêt puis au redémarrage du serveur.
\paragraph{\texttt{scripts/e2e-server.js}} Prépare et arrête l'environnement serveur/base isolé des tests navigateur.
\paragraph{\texttt{scripts/local-mongo.js}} Démarre MongoDB local persistant dans .local/mongodb avec WiredTiger.
\paragraph{\texttt{scripts/setup.js}} Initialise la configuration locale du backend sans remplacer un fichier déjà présent.
\section*{Périmètre et fichiers générés}
L'inventaire est basé sur les fichiers présents dans le dépôt au moment de la rédaction. Les dossiers \texttt{node\_modules/} contiennent les bibliothèques téléchargées par npm et changent selon l'installation; ils ne sont pas détaillés fichier par fichier. \texttt{frontend/dist/}, \texttt{test-results/}, les captures et prévisualisations sont générés par les commandes de build ou de test. \texttt{.local/mongodb/} est une vraie base persistante : ne pas supprimer ce dossier si l'on veut garder les comptes et tâches.

Pour les commandes d'inspection et les détails sur les collections MongoDB, voir \href{DONNEES.md}{docs/DONNEES.md}. Pour le contrat complet de l'API, voir \texttt{docs/openapi.json} ou ouvrir Swagger à \url{http://localhost:3000/api/docs/} après démarrage.
\end{document}
